%!TEX root = ../exa-ma-d7.3.tex
\section{Next Steps \& Roadmap for Software Training}
\label{sec:next-steps-software}

Over the next months, our primary goals for software-related training are threefold: 
\begin{itemize}
    \item \textbf{Refine and update existing training materials} to address early feedback and integrate lessons learned from the initial demonstration phase.
    \item \textbf{Develop new materials} to expand coverage of Exa-MA software, emphasizing novel functionalities and best practices for exascale architectures.
    \item \textbf{Document the Exa-MA software development process} to ensure consistent, reproducible, and sustainable practices across all Work Packages.
\end{itemize}

\subsection{Updating Existing Training Materials}

Currently, each Exa-MA software stack (e.g., \textit{FreeFem++}, \textit{Feel++}, \textit{HPDDM}, \textit{Mmg/Parmmg}, etc.) offers preliminary tutorials, user guides, or workshop slides (see Table~\ref{tab:software-training}). 
%
Over the next six months, we plan to:
\begin{itemize}
    \item \textbf{Incorporate user feedback} gained from the first round of training sessions and hands-on labs into these materials. 
    \item \textbf{Enrich examples} with advanced use cases (e.g., domain decomposition, GPU-acceleration tips, container-based workflows) to reflect ongoing progress in WP1--WP6.
    \item \textbf{Streamline access} by centralizing PDF slide decks, tutorial notebooks, and video content on a dedicated Exa-MA online portal.
\end{itemize}
These updates will ensure that current and future participants—ranging from early-career researchers to HPC experts—can more easily adopt the software.

\subsection{Adding New Training Material}

In parallel to refining existing content, we will create additional resources to address evolving needs across the Exa-MA ecosystem:
\begin{itemize}
    \item \textbf{Mini-App Tutorials:} 
    New, step-by-step guides will highlight mini- and proxy-app integrations within Exa-MA, emphasizing real-world examples and showcasing how to measure and optimize performance at scale.
    \item \textbf{Advanced Topics:}
    We plan to develop short courses on critical exascale issues such as resilience and fault tolerance (addressing B9), mixed precision methods for solvers (B7/B2), and large-scale data handling (B6).
    \item \textbf{Cross-WP Synergy:}
    The new materials will integrate insights from WP2 (model order reduction), WP3 (multi-physics solvers), and WP5 (large-scale optimization), ensuring consistent methodology and terminology across Exa-MA.
\end{itemize}

\subsection{Documenting the Exa-MA Software Development Process}

To guarantee reproducibility and maintain a high standard of code quality, a \textbf{development process handbook} will be introduced. This handbook will:
\begin{itemize}
    \item \textbf{Outline best practices} for version control, continuous integration, and containerization of Exa-MA software, aligned with the guidelines established in WP7.
    \item \textbf{Provide a roadmap} for contributors to follow when implementing new features, writing documentation, or preparing releases.
    \item \textbf{Detail our co-design approach} (see Task 7.2) by describing how software enhancements and performance benchmarks feed back into WP1--WP6 research objectives.
\end{itemize}

In addition, we will foster a more transparent development environment by encouraging all contributors to log progress and share milestone artifacts (e.g., documentation changes, test results) on a publicly accessible repository.

\subsection{Milestones and KPI Tracking}

By M24, we aim to have:
\begin{itemize}
    \item \textbf{At least 20\%} of existing tutorials upgraded with advanced HPC examples (GPU usage, multi-node scaling) in particular using Kokkos.
    \item \textbf{Two or more new training modules} focusing on mini-app usage in actual exascale scenarios.
    \item \textbf{A finalized development process handbook}, used by all contributors to ensure consistent practices across Exa-MA.
\end{itemize}
We will measure progress via:
\begin{itemize}
    \item \textbf{Attendance metrics} at upcoming training events (both in-person workshops and online webinars).
    \item \textbf{Repository analytics} (e.g., download counts, Git commits) for updated/new training content.
    \item \textbf{Feedback surveys} after each training session, capturing ease of adoption and clarity of the new materials.
\end{itemize}

